% GENERATED FILE. Edit canonical lessons, then run npm run generate:books.
% canonical-source-hash: fnv1a64:2174205b6cac96fb
% canonical-lessons: LA-C138-piscator, LA-C138-servus, LA-C138-domina, LA-C138-amica, LA-C138-armarium

\chapter{Fisherman, Slave, Mistress, Friend, Cupboard}
\label{ch:la-fisherman-slave-mistress-friend-cupboard}
\hlchaptermodality{138}

\begin{chapteropening}
\textbf{By the end of this chapter:} \emph{I can say a fisherman, a slave, a mistress, a friend and a cupboard.}

Complete the last lesson of chapter 138: I can say a fisherman, a slave, a mistress, a friend and a cupboard.
\end{chapteropening}

\section[piscātor, piscātōris]{piscātor, piscātōris — a fisherman}

\label{lesson:LA-C138-piscator}



\begin{quote}
Before the new one: say the Latin for a barber, then the Latin for a shepherd.
\end{quote}

\subsection*{You'll want to know: piscātor, piscātōris}
\textbf{piscātor, piscātōris} — \textquotedblleft{}a fisherman\textquotedblright{}.

\begin{cousinweb}
\textbf{piscātor} is built on \emph{piscor}, \textquotedblleft{}to fish\textquotedblright{}, from \emph{piscis}, \textquotedblleft{}a fish\textquotedblright{} (EtymDB 2.1), with \emph{-tor}, \textquotedblleft{}one who does it\textquotedblright{}. \emph{piscis} is usually traced to Proto-Indo-European \emph{*pisk-}, \textquotedblleft{}fish\textquotedblright{} (EtymDB 2.1); English \emph{fish} is a cousin.
\end{cousinweb}

\begin{grammarlens}[title={What you've built}]
One more word for this chapter.
\end{grammarlens}

\subsection*{Your turn}
\begin{itemize}
\raggedright
  \item \emph{Say it:} \emph{piscātor}
  \item \emph{Say it:} \emph{piscātor}, once more
  \item \emph{Say it:} say \emph{piscātor}
  \item \emph{Recall:} say the Latin for a barber, then the Latin for a shepherd, then say \emph{piscātor} again
\end{itemize}

\begin{tcolorbox}[breakable,colback=teal!4,colframe=teal!35!black,title={Before you move on}]
What does piscātor, piscātōris mean? (\textquotedblleft{}A fisherman\textquotedblright{}.) Say it once more.
\end{tcolorbox}
\section[servus, servī]{servus, servī — a slave, a servant}

\label{lesson:LA-C138-servus}



\begin{quote}
Before the new one: say the Latin for a shepherd, then the Latin for a fisherman.
\end{quote}

\subsection*{You'll want to know: servus, servī}
\textbf{servus, servī} — \textquotedblleft{}a slave, a servant\textquotedblright{}.

\begin{cousinweb}
\textbf{servus} has no sure origin. Romans tied it to \emph{servō}, \textquotedblleft{}to keep, to save\textquotedblright{} (Lewis \& Short), and an Etruscan source has also been suggested; neither is proven.
\end{cousinweb}

\begin{grammarlens}[title={What you've built}]
One more word for this chapter.
\end{grammarlens}

\subsection*{Your turn}
\begin{itemize}
\raggedright
  \item \emph{Say it:} \emph{servus}
  \item \emph{Say it:} \emph{servus}, once more
  \item \emph{Say it:} say \emph{servus}
  \item \emph{Recall:} say the Latin for a shepherd, then the Latin for a fisherman, then say \emph{servus} again
\end{itemize}

\begin{tcolorbox}[breakable,colback=teal!4,colframe=teal!35!black,title={Before you move on}]
What does servus, servī mean? (\textquotedblleft{}A slave, a servant\textquotedblright{}.) Say it once more.
\end{tcolorbox}
\section[domina, dominae]{domina, dominae — a mistress, the lady of the house}

\label{lesson:LA-C138-domina}



\begin{quote}
Before the new one: say the Latin for a fisherman, then the Latin for a slave.
\end{quote}

\subsection*{You'll want to know: domina, dominae}
\textbf{domina, dominae} — \textquotedblleft{}a mistress, the lady of the house\textquotedblright{}.

\begin{cousinweb}
\textbf{domina} is the feminine of \emph{dominus}, \textquotedblleft{}master of the house\textquotedblright{} (Lewis \& Short), built on \emph{domus}, \textquotedblleft{}house\textquotedblright{}, which is usually traced to Proto-Indo-European \emph{*dom-}, \textquotedblleft{}house\textquotedblright{} (EtymDB 2.1).
\end{cousinweb}

\begin{grammarlens}[title={What you've built}]
One more word for this chapter.
\end{grammarlens}

\subsection*{Your turn}
\begin{itemize}
\raggedright
  \item \emph{Say it:} \emph{domina}
  \item \emph{Say it:} \emph{domina}, once more
  \item \emph{Say it:} say \emph{domina}
  \item \emph{Recall:} say the Latin for a fisherman, then the Latin for a slave, then say \emph{domina} again
\end{itemize}

\begin{tcolorbox}[breakable,colback=teal!4,colframe=teal!35!black,title={Before you move on}]
What does domina, dominae mean? (\textquotedblleft{}A mistress, the lady of the house\textquotedblright{}.) Say it once more.
\end{tcolorbox}
\section[amīca, amīcae]{amīca, amīcae — a friend (a woman)}

\label{lesson:LA-C138-amica}



\begin{quote}
Before the new one: say the Latin for a slave, then the Latin for a mistress.
\end{quote}

\subsection*{You'll want to know: amīca, amīcae}
\textbf{amīca, amīcae} — \textquotedblleft{}a friend (a woman)\textquotedblright{}.

\begin{cousinweb}
\textbf{amīca} is the feminine of \emph{amīcus}, \textquotedblleft{}a friend\textquotedblright{} (Lewis \& Short), which is built on \emph{amō}, \textquotedblleft{}to love\textquotedblright{}.
\end{cousinweb}

\begin{grammarlens}[title={What you've built}]
One more word for this chapter.
\end{grammarlens}

\subsection*{Your turn}
\begin{itemize}
\raggedright
  \item \emph{Say it:} \emph{amīca}
  \item \emph{Say it:} \emph{amīca}, once more
  \item \emph{Say it:} say \emph{amīca}
  \item \emph{Recall:} say the Latin for a slave, then the Latin for a mistress, then say \emph{amīca} again
\end{itemize}

\begin{tcolorbox}[breakable,colback=teal!4,colframe=teal!35!black,title={Before you move on}]
What does amīca, amīcae mean? (\textquotedblleft{}A friend (a woman)\textquotedblright{}.) Say it once more.
\end{tcolorbox}
\section[armārium, armāriī]{armārium, armāriī — a cupboard}

\label{lesson:LA-C138-armarium}



\begin{quote}
Before the new one: say the Latin for a mistress, then the Latin for a friend.
\end{quote}

\subsection*{You'll want to know: armārium, armāriī}
\textbf{armārium, armāriī} — \textquotedblleft{}a cupboard\textquotedblright{}.

\begin{cousinweb}
\textbf{armārium} is built on \emph{arma}, \textquotedblleft{}tools, gear\textquotedblright{} (Lewis \& Short), with \emph{-ārium}, an ending for a place to keep things.
\end{cousinweb}

\begin{grammarlens}[title={What you've built}]
One more word for this chapter.
\end{grammarlens}

\subsection*{Your turn}
\begin{itemize}
\raggedright
  \item \emph{Say it:} \emph{armārium}
  \item \emph{Say it:} \emph{armārium}, once more
  \item \emph{Say it:} say \emph{armārium}
  \item \emph{Recall:} say the Latin for a mistress, then the Latin for a friend, then say \emph{armārium} again
\end{itemize}

\begin{tcolorbox}[breakable,colback=teal!4,colframe=teal!35!black,title={Before you move on}]
What does armārium, armāriī mean? (\textquotedblleft{}A cupboard\textquotedblright{}.) Say it once more.
\end{tcolorbox}
